% =============================================================================
%  CAGE — TECHNICAL COMPANION
%  A comprehensive academic companion to "CAGE: An Evaluation Framework for
%  Cache-Augmented Generation Models" (see updated-my-article.tex).
%
%  Scope: explains the whole project — inference internals, vLLM, RAG vs CAG
%  (how each works and how each fails), IR-quality / hallucination testing,
%  the cage-stats telemetry tool, KV/prompt compression, and the phase roadmap.
%
%  Every project-specific claim is grounded in the CAGE/cage-stats codebase;
%  every external claim is grounded in a cited, verified reference (see the
%  Document Validation section, which lists provenance and known caveats).
%
%  Preamble matches updated-my-article.tex so it compiles in the same
%  environment (SBC template). Bibliography is a manual thebibliography
%  (no .bst / BibTeX pass required).
% =============================================================================
\documentclass[12pt,a4paper]{article}
\usepackage{float}
\usepackage[section]{placeins}
\usepackage{sbc-template}
\usepackage[english]{babel}
\usepackage{times}
\usepackage{courier}
\usepackage[scaled]{helvet}
\usepackage{array}
\usepackage{enumitem}
\hbadness=10000
\vbadness=10000
\setlength{\emergencystretch}{3em}
\usepackage[utf8]{inputenc}
\usepackage[T1]{fontenc}
\usepackage{url}
\usepackage[authoryear]{natbib}
\bibpunct{[}{]}{;}{a}{ }{,}
\usepackage{graphicx}
\usepackage{amsmath}
\usepackage{hyperref}
\let\cite\citep
\usepackage{longtable}
\usepackage{booktabs}
\usepackage{setspace}
\usepackage{subcaption}
\usepackage{multirow}
\usepackage{xcolor}

\renewcommand{\topfraction}{0.9}
\renewcommand{\bottomfraction}{0.8}
\renewcommand{\textfraction}{0.07}
\renewcommand{\floatpagefraction}{0.8}

\title{CAGE: A Technical Companion --- Inference Internals, RAG versus Cache-Augmented Generation, Telemetry, and Compression on the Road to Distributed KV-Cache Evaluation under HPC}
\author{Lucas Mariano do Carmo\inst{1}, Wladmir Cardoso Brandão\inst{1}, Henrique Cota de Freitas\inst{1}}
\address{Pontifícia Universidade Católica de Minas Gerais (PUC Minas), Belo Horizonte, Brazil \email{lucas.mariano@sga.pucminas.br, wladmir@pucminas.br, cota@pucminas.br}}
\date{}

\begin{document}
\singlespacing
\maketitle

\begin{abstract}
This companion document complements the CAGE article by explaining, from the
ground up, the systems and evaluation concepts on which the framework rests, and
by documenting the developments, phases, and plans produced so far. It describes
how Large Language Model (LLM) inference works (three distinct framings), the
internal workings of the vLLM serving engine, and how Retrieval-Augmented
Generation (RAG) and Cache-Augmented Generation (CAG) each function --- and how
each may fail under hallucination, information-retrieval (IR) quality problems,
content-gathering errors, and High-Performance Computing (HPC) workloads. It
documents the \texttt{cage-stats} serving-telemetry tool built for the project,
and reviews recent (2024--2025) work on prompt and Key-Value (KV) cache
compression and its asymmetric impact on CAG versus RAG. The document closes with
the three-phase test plan and an explicit validation section that records the
provenance of every claim and flags known caveats in the preliminary Phase~1
results.
\end{abstract}

\section{Purpose and Scope}

The CAGE article introduces a framework that jointly evaluates serving
performance, cache and retrieval telemetry, and semantic quality for cache-aware
LLM systems. That article is necessarily compact. The purpose of this companion is
to give the full technical narrative behind it: why cache-aware serving is a
distinct evaluation problem, how the underlying machinery actually works, what new
tools and baselines were developed, and what the planned experiments are. It is
written to be self-contained and to be cross-checked against the project's source
code; Section~\ref{sec:validation} records the provenance of each claim.

Three commitments organise the text. First, every project-specific statement is
traceable to the CAGE or \texttt{cage-stats} codebase. Second, every external
technical claim is grounded in a cited, verifiable reference, with emphasis on
recent (2024--2025) literature. Third, the document is explicit about what is
\emph{measured}, what is \emph{simulated}, and what is \emph{planned}, because the
distinction is central to the project's scientific integrity.

Table~\ref{tab:acronyms} lists the acronyms used throughout.

\begin{table}[H]
\centering
\caption{Acronyms and abbreviations.}
\label{tab:acronyms}
\footnotesize
\begin{tabular}{@{}ll@{}}
\toprule
\textbf{Acronym} & \textbf{Expansion} \\
\midrule
CAG & Cache-Augmented Generation \\
CAGE & Cache-Augmented Generation Evaluation (this framework) \\
RAG & Retrieval-Augmented Generation \\
KV & Key-Value (cache) \\
LLM & Large Language Model \\
TTFT & Time-To-First-Token \\
TPOT & Time-Per-Output-Token \\
QPS / TPS & Queries / Tokens Per Second \\
NLI & Natural Language Inference \\
IR & Information Retrieval \\
HPC & High-Performance Computing \\
vLLM & Virtual LLM serving engine (PagedAttention) \\
FAISS & Facebook AI Similarity Search (vector index) \\
NCCL & NVIDIA Collective Communications Library \\
gVNIC & Google Virtual Network Interface Card \\
MLA & Multi-head Latent Attention \\
FP8 & 8-bit floating-point (KV quantisation) \\
MoE & Mixture-of-Experts \\
GQA / MHA & Grouped-Query / Multi-Head Attention \\
\bottomrule
\end{tabular}
\end{table}

\section{How LLM Inference Works (Three Framings)}
\label{sec:inference}

Autoregressive Transformer inference is the substrate of this entire project: the
quantity CAG reuses --- the KV cache --- exists only because of how inference is
structured. Because different readers find different framings clearest, this
section explains inference three ways, then states why it sets the experimental
design.

\subsection{Framing 1 --- The Two-Phase View (prefill then decode)}

A request to an LLM is served in two phases. In the \emph{prefill} phase, the
model reads the entire input prompt in one parallel pass, computing for every
token a set of key and value vectors at each attention layer; these are written
into the KV cache~\cite{vaswani2017attention, kwon2023efficient}. In the
\emph{decode} phase, the model generates output one token at a time; each new
token attends to all cached keys and values and appends its own. Prefill is
compute-heavy and parallel; decode is memory-bandwidth-bound and sequential. The
two phases have different costs and are measured by different metrics ---
Time-To-First-Token (TTFT) is dominated by prefill, while Time-Per-Output-Token
(TPOT) characterises decode. This separation is so fundamental that production
systems now disaggregate the two phases onto different
hardware~\cite{zhong2024distserve}.

\subsection{Framing 2 --- The Memory-Reuse View}

Equivalently, inference can be seen as an exercise in \emph{not recomputing}.
Without a cache, generating token $n$ would require recomputing attention over all
$n-1$ previous tokens, giving cost quadratic in sequence length. The KV cache
trades computation for memory: by storing the per-token keys and values once, each
decoding step reuses them in $O(n)$ work. The price is memory --- the cache grows
linearly with sequence length and batch size and frequently exceeds the size of
the model weights~\cite{li2024kvsurvey}. From this angle, an LLM server is a
memory-management system as much as a compute system, and the central engineering
question is \emph{which cached state is worth keeping, moving, or
recomputing}. CAG is precisely the decision to keep and reuse the KV state of a
known, stable context instead of recomputing it.

\subsection{Framing 3 --- The Economic / Throughput View}

A third framing treats inference as a throughput economy under a fixed memory
budget. A GPU can hold only so many KV-cache blocks; every block spent on one
request's history is unavailable to batch another request. Serving efficiency is
therefore the art of packing as many concurrent sequences as possible into the KV
budget while keeping latency acceptable. Reuse (prefix caching) increases
effective capacity by letting requests share blocks; compression increases it by
shrinking each block; eviction reclaims it by discarding low-value
blocks~\cite{kvcompresssurvey}. In this view, the metrics that matter ---
throughput (QPS, tokens/s), TTFT, and tail latency --- are all expressions of how
well the system monetises a scarce memory resource.

\subsection{Why This Sets the Experimental Design}

These framings explain the metric suite and the baselines. Because prefill and
decode differ (Framing~1), CAGE separates TTFT from TPOT rather than reporting only
end-to-end latency. Because the bottleneck is memory reuse (Framing~2), CAGE
instruments cache-hit and prompt-token-source signals, not just request timings.
Because efficiency is a budget economy (Framing~3), the project's later phases
deliberately induce memory pressure (larger models, longer contexts) to observe
eviction and degradation --- effects that simply do not appear on a small,
unsaturated local node. The systems substrate, in short, dictates what must be
measured.

\section{Inside vLLM: The Serving Substrate}
\label{sec:vllm}

CAGE builds on vLLM, whose design choices determine what cache-aware evaluation
can observe. This section documents the specific mechanisms the project uses or
plans to use; the flag strings below are those actually referenced in the
codebase.

\paragraph{PagedAttention and prefix caching.} vLLM stores the KV cache in
fixed-size, non-contiguous blocks (paging), which reduces fragmentation and lets
distinct requests \emph{share} the blocks of a common prompt
prefix~\cite{kwon2023efficient}. CAGE enables this with
\texttt{-{}-enable-prefix-caching}; it is the mechanism behind the
\texttt{prefix\_cache}, \texttt{hybrid}, and \texttt{distributed} baselines.

\paragraph{Prompt-token accounting.} With \texttt{-{}-enable-prompt-tokens-details},
vLLM reports how many prompt tokens were served from cache versus recomputed. CAGE
uses this to compute the prompt-cached ratio --- the signal that distinguishes
``retrieval succeeded'' from ``computation was actually reused''.

\paragraph{The metrics endpoint.} vLLM exposes a Prometheus \texttt{/metrics}
endpoint. CAGE's telemetry component (Section~\ref{sec:cagestats}) scrapes it to
recover signals not visible at the request layer: prefix-cache hit rate,
speculative-decode acceptance, and KV-cache dtype/ratio.

\paragraph{Speculative decoding.} A small draft model proposes several tokens that
the target model verifies in parallel, accelerating decode when acceptance is
high~\cite{leviathan2023speculative}. CAGE configures it through a server-side
\texttt{VLLM\_SPECULATIVE\_CONFIG} (e.g.\ an n-gram method) and records the
acceptance rate from \texttt{/metrics}; this backs the \texttt{speculative}
baseline.

\paragraph{KV-cache quantisation.} vLLM can store the cache in 8-bit floating point
via \texttt{-{}-kv-cache-dtype fp8}, roughly halving its footprint. CAGE records
this setting (passed to the server through \texttt{VLLM\_KV\_CACHE\_DTYPE}) for the
\texttt{compressed\_cag} baseline.

\paragraph{Disaggregated prefill and cross-node KV transfer.} vLLM supports
disaggregated prefilling, in which the KV cache produced by a prefill worker is
transferred to a separate decode worker, typically through a KV-transfer layer
such as LMCache~\cite{vllm2024disagg, lmcache}. This is the mechanism CAGE's
Phase~3 targets for \emph{real} cross-node KV-cache movement. It is important to
state plainly (and Section~\ref{sec:validation} repeats it) that the current
\texttt{distributed} baseline performs HTTP-level routing with a modelled transfer
cost, \emph{not} real tensor transfer; wiring \texttt{-{}-kv-transfer-config} with
LMCache/NIXL is planned, not yet measured.

\section{RAG versus CAG, Top to Bottom}
\label{sec:ragvscag}

This section explains both paradigms mechanistically and then catalogues how each
fails. The two are not simply ``retrieve'' versus ``cache''; they make opposite
bets about where knowledge should live at inference time.

\subsection{How RAG Works}

Retrieval-Augmented Generation grounds a model in an external corpus at query
time~\cite{lewis2020retrieval}. A typical pipeline: (i) embed the query with a
dense encoder; (ii) search a vector index (in CAGE, FAISS with exact inner-product
search) for nearest-neighbour passages; (iii) optionally rerank candidates with a
cross-encoder; (iv) concatenate the top-$k$ passages into the prompt; (v) generate.
RAG's strength is freshness and coverage: knowledge can be updated by editing the
corpus, with no model retraining. Its cost is paid \emph{every request} --- an
embedding pass, an index search, a rerank, and the processing of freshly inserted
(uncached) context tokens.

\subsection{How CAG Works}

Cache-Augmented Generation makes the opposite bet: when the relevant knowledge is
stable and small enough to fit the context window, load it once, precompute its KV
cache, and reuse that cache for every subsequent query, computing only the new
query tokens~\cite{yu2024dontdorag}. There is no retrieval step at inference time;
the ``retrieval'' happened once, offline, when the context was chosen and cached.
CAG's strength is latency and simplicity: it removes the retrieval pipeline and
turns prompt processing into a near-free cache lookup. Its cost is paid \emph{up
front and in memory} --- the context must fit, the cache must be held, and updating
the knowledge means recomputing the cache.

\subsection{The Core Differences}

Table~\ref{tab:ragvscag} summarises the contrast. The decisive variables are
\emph{where knowledge lives} (external corpus vs.\ resident cache), \emph{when the
cost is paid} (per-request vs.\ up-front), and \emph{what limits scale} (retrieval
quality and pipeline latency vs.\ context window and KV memory).

\begin{table}[H]
\centering
\caption{RAG versus CAG along the dimensions that matter for serving and quality.}
\label{tab:ragvscag}
\footnotesize
\begin{tabular}{@{}p{3.2cm}p{4.6cm}p{4.6cm}@{}}
\toprule
\textbf{Dimension} & \textbf{RAG} & \textbf{CAG} \\
\midrule
Knowledge location & External corpus, fetched per query & Resident in context / KV cache \\
When cost is paid & Every request (embed + search + rerank) & Once, up front (preload + precompute KV) \\
Dominant latency & Retrieval + uncached prefill & Near-zero prefill on cache hit; decode only \\
Knowledge freshness & High (edit the corpus) & Low (recompute cache to update) \\
Scale limit & Corpus size is unbounded; quality bounded by retriever & Bounded by context window and KV memory \\
Primary failure mode & Retrieval error / IR quality & Staleness, context overflow, KV pressure \\
\bottomrule
\end{tabular}
\end{table}

Because neither paradigm dominates universally, recent work frames the choice as a
\emph{routing} problem --- selecting retrieval or long-context/cached processing per
query based on the model's own assessment of need~\cite{li2024selfroute} --- a
direction that CAGE's hybrid baselines are designed to evaluate.

\subsection{How RAG May Fail}
\label{sec:ragfail}

\paragraph{Retrieval error and content gathering.} RAG inherits every weakness of
its retriever. If the index returns an irrelevant, partial, or wrong passage, the
generator is grounded in bad evidence. Barnett et al.~\cite{phillips2024seven}
catalogue seven recurring engineering failure points of production RAG --- among
them missing content, passages not retrieved, passages not in context after
reranking, and extraction failures --- showing that ``content gathering'' is a
first-class reliability problem, not a solved preprocessing step.

\paragraph{Hallucination despite retrieval.} Crucially, retrieving the right
passage does not guarantee a grounded answer: the model may still produce claims
unsupported by or contradicting the retrieved text. RAGTruth~\cite{ragtruth}
documents this at the span level with $18{,}000$ annotated examples, and is the
training corpus for modern detectors. When the gold context is \emph{already}
available, retrieval can even \emph{reduce} grounding by injecting topically
similar but imperfect passages~\cite{yu2024dontdorag} --- exactly the effect the
CAGE Phase~1 quality comparison was designed to expose. Broader surveys of
hallucination in foundation models frame this as an intrinsic, not incidental,
risk~\cite{zhang2024siren}.

\paragraph{IR-quality measurement is itself hard.} Because a high retrieval
hit-rate can coexist with poor answers, RAG must be evaluated with metrics that
test \emph{grounding}, not surface similarity. This motivated reference-free
evaluators such as RAGAS~\cite{espejel2023ragas} and NLI-judge frameworks such as
ARES~\cite{ares2024}, benchmark suites such as RAGBench~\cite{li2024ragbench} and
BERGEN~\cite{bergen2024}, and the IR-quality instrumentation in CAGE
(Section~\ref{sec:irquality}).

\paragraph{HPC and serving overheads.} At scale, the retrieval pipeline becomes a
serving liability: each request pays embedding, search, rerank, and the prefill of
\emph{uncached} context tokens. CAGE's Phase~1 data illustrate the latency penalty
of this per-request cost; the systems-serving literature treats retrieval-time work
as a goodput tax that competes with generation for the same
resources~\cite{inferenceserving}.

\subsection{How CAG May Fail}
\label{sec:cagfail}

\paragraph{Context-window and quadratic limits.} CAG only works if the knowledge
fits the context window. Beyond that, it cannot be applied, and even within it,
prefill cost grows with context length~\cite{yu2024dontdorag}. CAG is thus suited
to \emph{stable, bounded} knowledge, not open-domain corpora.

\paragraph{Staleness and update cost.} Because the cache encodes a fixed context,
any change to the knowledge requires recomputing the KV cache. For volatile data,
this update cost can erase CAG's latency advantage --- the mirror image of RAG's
freshness strength.

\paragraph{KV memory pressure and eviction under HPC.} CAG's resident cache is
exactly the object that dominates GPU memory at scale. Under HPC workloads --- many
concurrent long-context sessions --- the cache saturates VRAM, forcing eviction
policies to trigger and degrading throughput; this is the regime that motivates the
entire KV-compression literature~\cite{li2024kvsurvey, kvcompresssurvey}. A small,
unsaturated node never exhibits this, which is why CAGE's Phase~2/3 deliberately
provokes it.

\paragraph{Cold-start tails and cross-node transfer.} In a distributed deployment,
a replica whose local cache is cold pays full prefill, producing severe tail
latency; and if the cache must be \emph{moved} between nodes (disaggregated
prefill), the transfer has a real bandwidth and latency
cost~\cite{qin2024mooncake, lmcache}. The central question of the project's later
phases is precisely when reusing or moving a remote KV cache is cheaper than
retrieving and recomputing context locally --- the distributed efficiency-versus-IR-quality
trade-off.

\section{Information-Retrieval Quality and Hallucination Testing}
\label{sec:irquality}

Because both paradigms can be efficient yet wrong, CAGE treats answer quality as a
first-class, instrumented axis. The metrics below are those implemented in the
project's evaluation module; their model identifiers are recorded here for
reproducibility.

\paragraph{Grounding (primary).} The primary quality signal is span-level
hallucination detection via LettuceDetect~\cite{lettucedetect}, a token-classification
model built on the long-context ModernBERT encoder~\cite{modernbert} and trained on
RAGTruth~\cite{ragtruth}. It processes \emph{(context, question, answer)} triples
and localises unsupported spans, yielding a grounding score and a
hallucinated-span ratio. The project uses the
\texttt{KRLabsOrg/lettucedect-base-modernbert-en-v1} model.

\paragraph{Faithfulness (NLI).} As a complementary and fallback signal, CAGE
computes a strict claim-level NLI faithfulness score: the answer is split into
claims, and each claim's maximum entailment probability against the context is
taken using \texttt{MoritzLaurer/DeBERTa-v3-base-mnli-fever-anli} (with
\texttt{facebook/bart-large-mnli} as fallback). The entailment label index is
resolved from the model configuration rather than hard-coded.

\paragraph{Negative control.} BERTScore (rescaled with baseline) is retained
\emph{not} as a primary metric but as a deliberate negative control: it measures
surface and vocabulary similarity, and the project's data show it is largely
non-discriminative across cache and retrieval settings --- empirically reinforcing
the argument of RAGAS~\cite{espejel2023ragas} and ARES~\cite{ares2024} that
soft-matching metrics miss factual hallucination.

\paragraph{Retriever diagnostics and QA standards.} Context relevance
(question-context cosine similarity) is reported as a retriever diagnostic, not as
answer quality, and standard SQuAD-style token F1 / Exact Match are included for
comparability. A key engineering correction recorded in the project is that all
quality metrics return a null value (excluded from means) when their model is
unavailable, rather than a sentinel constant --- the fix for an artefact that
affected earlier numbers (Section~\ref{sec:validation}).

\section{cage-stats: Serving Telemetry for CAGE}
\label{sec:cagestats}

A request-level harness sees latency and tokens, but not \emph{why} a request was
fast or slow. To close that gap, the project developed \texttt{cage-stats}, a
headless serving-telemetry tool (conceptually an ``nvtop for vLLM'') that polls the
engine and exposes the internal state CAGE needs.

\paragraph{What it captures.} \texttt{cage-stats} renders panels for concurrency
(running/waiting requests), throughput (generation and prompt tokens/s), latency
quantiles (p50/p95/p99), the KV-cache panel (compression ratio, dtype,
prefix-hit rate), session state, speculative-decoding acceptance and draft
tokens/step, an inference-efficiency figure, and multi-vendor GPU statistics (VRAM,
utilisation, power, temperature, PCIe bandwidth). The speculative-acceptance and
KV-ratio panels are exactly the signals that the request layer cannot expose.

\paragraph{How it integrates.} CAGE consumes \texttt{cage-stats} through a thin
bridge with a three-tier resolution order so the framework never hard-fails if the
tool is absent: (i) an in-process import of the \texttt{cage\_stats} headless API
(richest, honouring a \texttt{CAGE\_STATS\_HOME} path); (ii) a CLI subprocess
(\texttt{cage-stats -{}-once -{}-json}); (iii) graceful skip returning null. A
\texttt{mock} mode supplies synthetic data so dry-runs and tests need no live
server. The \texttt{-{}-vllm-telemetry} flag on the experiment runner captures a
snapshot at the end of each baseline, prints a one-shot dashboard, and persists a
\texttt{vllm\_telemetry.json} artefact alongside the request-level results, so cloud
runs record serving internals automatically.

\paragraph{Why it matters for the thesis.} The Phase~1 finding that \emph{local
reuse}, not retrieval success, drives TTFT is a telemetry claim: it can only be
made because the prompt-token-source and cache-hit signals are recorded.
\texttt{cage-stats} generalises this instrumentation to the GPU and distributed
regimes, where the new signals (KV compression ratio, cross-node transfer bytes,
speculative acceptance) become the evidence for the distributed trade-off study.

\section{Compression and Its Impact on CAG and RAG}
\label{sec:compression}

Compression is where the CAG/RAG distinction becomes sharp, because the two
paradigms compress \emph{different objects}. CAGE models this as an explicit axis
crossed with the context source, giving a 2$\times$2 design (full/compressed
$\times$ CAG/RAG).

\subsection{Two Families of Compression}

\paragraph{Prompt (token) compression --- helps RAG.} Prompt-compression methods
prune or rewrite the \emph{text} of the context before it reaches the model.
LLMLingua compresses prompts by removing low-information
tokens~\cite{llmlingua}; LLMLingua-2 reframes this as a token-classification task
distilled from a stronger model, achieving 2--5$\times$ compression with strong
generalisation~\cite{llmlingua2}. Because RAG inserts fresh, uncached passages on
every request, shrinking those passages directly reduces the (uncached) prefill it
must pay --- so prompt compression is naturally a \emph{RAG} lever. CAGE implements
this as the \texttt{compressed\_rag} baseline (LLMLingua-2,
\texttt{-{}-compress-method}, \texttt{-{}-compress-ratio}).

\paragraph{KV-cache compression --- helps CAG.} KV-cache methods compress the
\emph{cached tensors} rather than the text. Recent surveys organise the field into
quantisation (e.g.\ FP8 and mixed-precision schemes such as
KVTuner~\cite{kvtuner}), eviction (selecting which tokens to keep), and low-rank /
novel-attention methods~\cite{kvcompresssurvey, li2024kvsurvey,
javidnia2025kvccompress}. The most architectural of these is Multi-head Latent
Attention (MLA), introduced in DeepSeek-V2, which compresses the KV cache into a
low-rank latent vector and reduces KV memory by over 90\%~\cite{deepseekv2}.
Because CAG's cost \emph{is} the resident KV cache, shrinking that cache is a
\emph{CAG} lever. CAGE implements this as the \texttt{compressed\_cag} baseline
(\texttt{-{}-kv-cache-dtype fp8}, with MLA as an architectural variant).

\subsection{What CAGE Runs (Implemented Stack)}

Table~\ref{tab:compstack} records the concrete libraries, algorithms, and models
behind each compression arm, for reproducibility. The prompt-compression arm is a
client-side dependency (the \texttt{llmlingua} package, MIT-licensed), so it runs on
any node; the KV-compression arm is server-side and native to vLLM (FP8) or
architectural (MLA), so it requires GPU execution. If \texttt{llmlingua} is absent,
\texttt{compressed\_rag} degrades gracefully to no compression (ratio $1.0$) rather
than failing.

\begin{table}[H]
\centering
\caption{Implemented compression stack: libraries, algorithms, and models.}
\label{tab:compstack}
\footnotesize
\begin{tabular}{@{}p{2.6cm}p{2.2cm}p{3.2cm}p{3.4cm}@{}}
\toprule
\textbf{Arm} & \textbf{Library} & \textbf{Algorithm} & \textbf{Model / setting} \\
\midrule
\texttt{compressed\_rag} (default) & \texttt{llmlingua}~\cite{llmlingua2} & LLMLingua-2 token classification (keep/drop) & \texttt{llmlingua-2-xlm-roberta-large-meetingbank}; keep ratio $0.5$ ($\sim$2$\times$) \\
\texttt{compressed\_rag} (alt) & \texttt{llmlingua}~\cite{llmlingua} & LLMLingua v1 (perplexity-based) & Llama-2-7B causal LM (heavier) \\
\texttt{compressed\_cag} (FP8) & vLLM (native) & FP8 KV quantisation~\cite{kvtuner} & \texttt{-{}-kv-cache-dtype fp8} at server launch \\
\texttt{compressed\_cag} (MLA) & model architecture & Multi-head Latent Attention~\cite{deepseekv2} & DeepSeek-V2-class model \\
\bottomrule
\end{tabular}
\end{table}

\subsection{The 2$\times$2 Design and Why the Asymmetry Matters}

Table~\ref{tab:compression} crosses context source with compression. The asymmetry
is the point: prompt compression attacks RAG's per-request token cost, while KV
compression attacks CAG's resident-memory cost. They are not interchangeable, and a
fair evaluation must apply each to the paradigm it actually helps. CAGE therefore
reports compression-specific metrics --- the token-retention ratio for prompt
compression, and KV bytes-per-token (and, in Phase~3, cross-node transfer bytes)
for KV compression --- so that an efficiency gain can always be weighed against any
quality loss it causes.

\begin{table}[H]
\centering
\caption{The compression axis crossed with the context source. Each cell names the
CAGE baseline and the object being compressed.}
\label{tab:compression}
\footnotesize
\begin{tabular}{@{}lll@{}}
\toprule
 & \textbf{CAG (cached context)} & \textbf{RAG (retrieved context)} \\
\midrule
\textbf{Full} & \texttt{prefix\_cache} --- full KV & \texttt{rag} --- full passages \\
\textbf{Compressed} & \texttt{compressed\_cag} --- FP8 / MLA KV & \texttt{compressed\_rag} --- LLMLingua-2 text \\
\bottomrule
\end{tabular}
\end{table}

\paragraph{Fairness: match the ratio, not the algorithm.} A natural objection is
that the two arms use different algorithms, so the comparison is unfair. The
objection mistakes the question: CAGE does not compare compression \emph{algorithms}
(LLMLingua-2 versus FP8) --- it compares serving \emph{paradigms}, each under its
own native compression, because text tokens and KV tensors are different objects to
which no single algorithm applies. Fairness is instead enforced by three controls
that \emph{are} identical across arms: the same quality metrics, the same protocol
(model, dataset, seeds, trials, decoding), and a \emph{matched compression ratio}.
Concretely, the headline KV arm is pinned to FP8 ($\sim$2$\times$ bytes per token)
so that it is ratio-matched to LLMLingua-2 at keep-ratio $0.5$ ($\sim$2$\times$
tokens); MLA ($\sim$7--14$\times$) is reported separately as a more-aggressive
architectural point, not as the matched comparator. Finally, the design admits no
contaminated comparison: Table~\ref{tab:compression} is read \emph{down} the columns
(CAG versus RAG, full or compressed) or \emph{across} the rows (the effect of
compression within one paradigm), never on the diagonal (algorithm versus
algorithm).

\subsection{Impact on CAG Architecture at Scale}

KV compression changes what is feasible for CAG. By reducing bytes-per-token, FP8
and especially MLA let a larger or longer context remain resident, push back the
saturation point at which eviction triggers, and --- decisively for Phase~3 --- cut
the volume that must cross the interconnect when a cache is transferred between
nodes~\cite{deepseekv2, lmcache}. In other words, compression is not a side topic:
it directly governs where CAG stops being viable and how expensive distributed CAG
is. This is why the compression axis is built into the framework rather than left
as future work.

\subsection{Pre-Registered Expectations}

Because the KV-compression arm requires GPU execution, its quantitative results are
produced in the GPU and HPC phases. The expectations in Table~\ref{tab:compexpect}
are therefore stated in advance, grounded in the cited literature, so that the
measured outcomes can be compared against a pre-registered hypothesis rather than
rationalised afterward.

\begin{table}[H]
\centering
\caption{Pre-registered expectations for the compression axis (to be confirmed or
refuted under the GPU/HPC protocol).}
\label{tab:compexpect}
\footnotesize
\begin{tabular}{@{}p{3.0cm}p{8.6cm}@{}}
\toprule
\textbf{Comparison} & \textbf{Expected outcome (hypothesis)} \\
\midrule
\texttt{compressed\_rag} vs \texttt{rag} & About 2$\times$ shorter prompts, lowering TTFT, latency, and cost at a small faithfulness loss (LLMLingua-2 reports near-lossless 2--5$\times$~\cite{llmlingua2}); it does \emph{not} close RAG's grounding gap. \\
\texttt{compressed\_cag} (FP8) vs \texttt{prefix\_cache} & About half the KV bytes per token at near-parity quality and latency~\cite{kvtuner}; the eviction threshold is pushed back. \\
\texttt{compressed\_cag} (MLA) & Over 90\% KV reduction~\cite{deepseekv2}; the largest lever for distributed CAG, since less cache must cross the interconnect in Phase~3. \\
Overall frontier & Compression shifts every baseline toward lower cost without re-ordering the CAG~$>$~RAG grounding ranking --- i.e.\ it is largely \emph{orthogonal} to the CAG-versus-RAG quality gap. \\
\bottomrule
\end{tabular}
\end{table}

The practical implication is that the choice of \emph{paradigm} (cache versus
retrieve) is expected to dominate answer quality, while compression is the knob for
cost and memory \emph{within} each paradigm --- and KV compression specifically is
what is expected to make distributed CAG affordable at HPC scale.

\section{Phases, Developments, and Test Plans}
\label{sec:phases}

The project advances through three phases that increase realism while preserving
measurement integrity. The baselines exercised are \texttt{no\_cache},
\texttt{prefix\_cache}, \texttt{rag}, \texttt{redis}, \texttt{hybrid},
\texttt{distributed}, \texttt{speculative}, \texttt{compressed\_rag}, and
\texttt{compressed\_cag}; datasets include SQuAD~v2~\cite{rajpurkar2018squad},
Natural Questions~\cite{kwiatkowski2019nq}, and the multi-hop
MuSiQue~\cite{trivedi2022musique}.

\paragraph{Phase 1 --- Local validation (completed).} A single CPU node
(Apple~M4~Pro, vLLM CPU backend) running Qwen3-4B on SQuAD~v2, seven baselines,
50 measured queries $\times$ 3 trials. The goal was structural and statistical
integrity of the orchestration and metric pipeline, not generalisable absolute
numbers. Its results, and their caveats, are discussed in
Section~\ref{sec:validation}.

\paragraph{Phase 2 --- GPU validation and scaling (next).} A single cloud GPU VM
(NVIDIA L4 on \texttt{g2-standard-8}; NVIDIA A100 on \texttt{a2-highgpu-1g} for
memory-pressure runs) re-runs the baselines on larger models (Qwen3-8B/14B) at
higher replication (e.g.\ 100 queries $\times$ 10 trials), under the corrected
protocol. This is where the compression axis is exercised under genuine memory
pressure and where the Phase~1 quality numbers are re-established with the fixed
metrics. Serving telemetry is captured automatically (Section~\ref{sec:cagestats}).

\paragraph{Phase 3 --- Distributed / HPC (future).} A multi-node A100 cluster (a
router plus replicas, with gVNIC and tuned NCCL topologies) evaluates \emph{real}
cross-node KV-cache transfer via disaggregated prefilling and
LMCache~\cite{vllm2024disagg, lmcache}, together with speculative-decoding
interactions. This phase quantifies the distributed efficiency-versus-IR-quality
frontier: when moving or reusing a remote KV cache beats retrieving and recomputing
context locally.

\paragraph{Test-plan structure.} Each phase reports the same triad --- serving
metrics (latency, TTFT, TPOT, throughput, tail percentiles), cache/transfer
telemetry, and quality (grounding, faithfulness, with BERTScore as control) ---
under a common protocol: fixed seed with per-trial reseeded sampling, cold-start
control via cache flushing where required, warmup excluded from measured metrics,
and non-parametric significance testing (Wilcoxon with Holm correction) with tail
percentiles reported alongside means.

\section{Document Validation and Provenance}
\label{sec:validation}

In keeping with the project's emphasis on avoiding unsupported claims, this section
records how each class of statement in this document was checked, and flags the
known caveats that any reader (or reviewer) must keep in mind.

\paragraph{Provenance policy.} Project-specific claims (baselines, flags, metric
models, telemetry behaviour, phase hardware) were taken from the CAGE and
\texttt{cage-stats} source and documentation, not from memory. External technical
claims were grounded in the cited references, each verified to exist with the
stated venue or arXiv identifier (e.g.\ LLMLingua-2, \emph{Findings of ACL 2024},
arXiv:2403.12968; LettuceDetect, arXiv:2502.17125; RAGTruth, \emph{ACL 2024},
arXiv:2401.00396; DeepSeek-V2/MLA, arXiv:2405.04434; LMCache, arXiv:2510.09665;
KV-compression survey, arXiv:2503.24000). Forward-looking statements are labelled
as plans.

\paragraph{Caveats on the preliminary Phase~1 numbers.} The absolute quality
constants reported in the Phase~1 validation (for instance the faithfulness and
BERTScore figures in the article) are \emph{preliminary} and were produced under an
evaluation pipeline that has since been corrected: the faithfulness and BERTScore
implementations were rewritten, LettuceDetect was adopted as the primary grounding
signal, the e5 retriever prefixes were fixed, and per-trial reseeding was added.
These specific constants should therefore be read as indicative of the framework's
behaviour, not as final measurements; Phase~2 re-establishes them under the
corrected protocol. The qualitative findings --- that local cache reuse drives
TTFT, that retrieval can reduce grounding when gold context is available, and that
soft-matching similarity is a poor hallucination detector --- are the robust claims
and are consistent with the cited literature.

\paragraph{The ``distributed'' caveat.} The Phase~1 \texttt{distributed} baseline
performs orchestration-level routing with a \emph{modelled} transfer cost; it does
\emph{not} move real KV tensors between nodes. Real cross-node KV transfer
(LMCache/NIXL, disaggregated prefill) is a Phase~3 plan. This document and the
framework state this explicitly so that ``distributed'' is never mistaken for a
measured distributed-systems result before Phase~3.

\paragraph{Known confounds to resolve in Phase~2/3.} Two protocol confounds are
documented and slated for correction: a gold-versus-retrieved context confound
(CAG arms receive the oracle passage while RAG receives retrieved passages), to be
resolved with an explicit context-source control; and warm-hybrid leakage (warmup
drawn from the measured set), to be resolved by disjoint warmup. Disclosing these
is part of the validation discipline, not a footnote.

\paragraph{Figures.} The figures in this document (Figures~\ref{fig:latency}--\ref{fig:pareto})
are the Phase~1 serving and quality plots. The serving-side figures (latency,
throughput, breakdown, Pareto) are systems measurements and are robust to the
metric corrections above; the quality figure is shown as preliminary, subject to
the caveats in this section.

\begin{figure}[H]
\centering
\includegraphics[width=0.8\textwidth]{companion_images/01_latency_comparison.png}
\caption{Phase~1 average latency and TTFT across baselines (preliminary, CPU node).}
\label{fig:latency}
\end{figure}

\begin{figure}[H]
\centering
\includegraphics[width=0.8\textwidth]{companion_images/04_latency_breakdown.png}
\caption{Phase~1 latency decomposition into prefill (TTFT) and generation. The
prefix-cache advantage is concentrated in prefill --- the systems claim that
survives metric correction.}
\label{fig:breakdown}
\end{figure}

\begin{figure}[H]
\centering
\includegraphics[width=0.8\textwidth]{companion_images/02_throughput_comparison.png}
\caption{Phase~1 throughput across baselines (queries/s and tokens/s).}
\label{fig:throughput}
\end{figure}

\begin{figure}[H]
\centering
\includegraphics[width=0.8\textwidth]{companion_images/03_quality_comparison.png}
\caption{Phase~1 quality metrics across baselines (preliminary; subject to the
metric-correction caveats in Section~\ref{sec:validation}).}
\label{fig:quality}
\end{figure}

\begin{figure}[H]
\centering
\includegraphics[width=0.8\textwidth]{companion_images/05_pareto_latency_quality.png}
\caption{Phase~1 latency-versus-quality plane, illustrating the efficiency-quality
trade-off that the distributed phases will quantify under HPC.}
\label{fig:pareto}
\end{figure}

\section{Conclusion}

This companion has explained the CAGE project from its systems foundations upward:
how inference produces and reuses the KV cache, how vLLM exposes the mechanisms
CAGE measures, how RAG and CAG each work and each fail, how IR quality and
hallucination are tested, how the \texttt{cage-stats} tool instruments the serving
layer, and how prompt and KV compression act asymmetrically on the two paradigms.
The through-line is the project's central trade-off --- distributed KV-cache
efficiency against information-retrieval quality --- and the discipline of stating
exactly what is measured, simulated, or planned at each phase. With the metric and
telemetry pipeline validated locally, the corrected protocol and the compression
axis carry the work into the GPU and HPC regimes where the trade-off becomes
quantifiable.

\begin{thebibliography}{99}

\bibitem[Barnett et~al.(2024)]{phillips2024seven}
Barnett, S., Kurniawan, S., Thudumu, S., Brannelly, Z., and Abdelrazek, M. Seven Failure Points When Engineering a Retrieval Augmented Generation System. In \textit{CAIN 2024}, pages 194--199, 2024.

\bibitem[Chan et~al.(2025)]{yu2024dontdorag}
Chan, B. J., Chen, C.-T., Cheng, J.-H., and Huang, H.-H. Don't Do RAG: When Cache-Augmented Generation is All You Need for Knowledge Tasks. In \textit{WWW (Companion Volume) 2025}, pages 893--897, 2025. arXiv:2412.15605.

\bibitem[DeepSeek-AI(2024)]{deepseekv2}
DeepSeek-AI. DeepSeek-V2: A Strong, Economical, and Efficient Mixture-of-Experts Language Model. \textit{arXiv preprint} arXiv:2405.04434, 2024.

\bibitem[Es et~al.(2024)]{espejel2023ragas}
Es, S., James, J., Espinosa Anke, L., and Schockaert, S. RAGAs: Automated Evaluation of Retrieval Augmented Generation. In \textit{Proceedings of EACL 2024: System Demonstrations}, pages 150--158, 2024.

\bibitem[Jiang et~al.(2023)]{llmlingua}
Jiang, H., Wu, Q., Lin, C.-Y., Yang, Y., and Qiu, L. LLMLingua: Compressing Prompts for Accelerated Inference of Large Language Models. In \textit{Proceedings of EMNLP 2023}, 2023. arXiv:2310.05736.

\bibitem[Javidnia et~al.(2025)]{javidnia2025kvccompress}
Javidnia, N., Rouhani, B. D., and Koushanfar, F. Key, Value, Compress: A Systematic Exploration of KV Cache Compression Techniques. In \textit{IEEE CICC 2025}, 2025.

\bibitem[Kovács et~al.(2025)]{lettucedetect}
Kovács, Á., Ország, G. M., Gémes, K., and Recski, G. LettuceDetect: A Hallucination Detection Framework for RAG Applications. \textit{arXiv preprint} arXiv:2502.17125, 2025.

\bibitem[Kwiatkowski et~al.(2019)]{kwiatkowski2019nq}
Kwiatkowski, T. et al. Natural Questions: A Benchmark for Question Answering Research. \textit{Transactions of the ACL}, 7:453--466, 2019.

\bibitem[Kwon et~al.(2023)]{kwon2023efficient}
Kwon, W., Li, Z., Zhuang, S., Sheng, Y., Zheng, L., Yu, C. H., Gonzalez, J. E., Zhang, H., and Stoica, I. Efficient Memory Management for Large Language Model Serving with PagedAttention. In \textit{Proceedings of the 29th SOSP}, pages 611--626, 2023.

\bibitem[Leviathan et~al.(2023)]{leviathan2023speculative}
Leviathan, Y., Kalman, M., and Matias, Y. Fast Inference from Transformers via Speculative Decoding. In \textit{Proceedings of ICML 2023}, 2023.

\bibitem[Lewis et~al.(2020)]{lewis2020retrieval}
Lewis, P. et al. Retrieval-Augmented Generation for Knowledge-Intensive NLP Tasks. In \textit{Advances in Neural Information Processing Systems (NeurIPS)}, 2020.

\bibitem[H. Li et~al.(2024)]{li2024kvsurvey}
Li, H. et al. A Survey on Large Language Model Acceleration based on KV Cache Management. \textit{arXiv preprint} arXiv:2412.19442, 2024.

\bibitem[Z. Li et~al.(2024)]{li2024selfroute}
Li, Z., Li, C., Zhang, M., Mei, Q., and Bendersky, M. Retrieval Augmented Generation or Long-Context LLMs? A Comprehensive Study and Hybrid Approach. In \textit{Proceedings of EMNLP 2024: Industry Track}, pages 881--893, 2024.

\bibitem[Friel et~al.(2024)]{li2024ragbench}
Friel, R. et al. RAGBench: Explainable Benchmark for Retrieval-Augmented Generation Systems. \textit{arXiv preprint} arXiv:2407.11005, 2024.

\bibitem[Niu et~al.(2024)]{ragtruth}
Niu, C., Wu, Y., Zhu, J., Xu, S., Shum, K., Zhong, R., Song, J., and Zhang, T. RAGTruth: A Hallucination Corpus for Developing Trustworthy Retrieval-Augmented Language Models. In \textit{Proceedings of ACL 2024}, 2024. arXiv:2401.00396.

\bibitem[Pan et~al.(2024)]{llmlingua2}
Pan, Z., Wu, Q., Jiang, H., Xia, M., Luo, X., Zhang, J., Lin, Q., Rühle, V., Yang, Y., Lin, C.-Y., Zhao, H. V., Qiu, L., and Zhang, D. LLMLingua-2: Data Distillation for Efficient and Faithful Task-Agnostic Prompt Compression. In \textit{Findings of the ACL 2024}, 2024. arXiv:2403.12968.

\bibitem[Qin et~al.(2024)]{qin2024mooncake}
Qin, R., Li, Z., He, W., Zhang, M., Wu, Y., Zheng, W., and Xu, X. Mooncake: A KVCache-centric Disaggregated Architecture for LLM Serving. \textit{arXiv preprint} arXiv:2407.00079, 2024.

\bibitem[Rajpurkar et~al.(2018)]{rajpurkar2018squad}
Rajpurkar, P., Jia, R., and Liang, P. Know What You Don't Know: Unanswerable Questions for SQuAD. In \textit{Proceedings of ACL 2018}, pages 784--789, 2018.

\bibitem[Rau et~al.(2024)]{bergen2024}
Rau, D., Déjean, H., Chirkova, N., Formal, T., Wang, S., Clinchant, S., and Nikoulina, V. BERGEN: A Benchmarking Library for Retrieval-Augmented Generation. In \textit{Findings of EMNLP 2024}, pages 7640--7663, 2024.

\bibitem[Saad-Falcon et~al.(2024)]{ares2024}
Saad-Falcon, J., Khattab, O., Potts, C., and Zaharia, M. ARES: An Automated Evaluation Framework for Retrieval-Augmented Generation Systems. In \textit{Proceedings of NAACL 2024 (Volume 1)}, pages 338--354, 2024.

\bibitem[Sahoo et~al.(2024)]{zhang2024siren}
Sahoo, P., Meharia, P., Ghosh, A., Saha, S., Jain, V., and Chadha, A. A Comprehensive Survey of Hallucination in Large Language, Image, Video and Audio Foundation Models. In \textit{Findings of EMNLP 2024}, pages 11709--11724, 2024.

\bibitem[Vaswani et~al.(2017)]{vaswani2017attention}
Vaswani, A., Shazeer, N., Parmar, N., Uszkoreit, J., Jones, L., Gomez, A. N., Kaiser, Ł., and Polosukhin, I. Attention Is All You Need. In \textit{Advances in Neural Information Processing Systems (NeurIPS)}, 2017.

\bibitem[vLLM Project(2024)]{vllm2024disagg}
vLLM Project. Disaggregated Prefilling (experimental). \textit{vLLM Documentation}, 2024. Available at: \url{https://docs.vllm.ai/en/stable/features/disagg_prefill.html}.

\bibitem[Warner et~al.(2024)]{modernbert}
Warner, B., Chaffin, A., Clavié, B., Weller, O., Hallström, O., Taghadouini, S., Gallagher, A., Biswas, R., Ladhak, F., Aarsen, T., Cooper, N., Adams, G., Howard, J., and Poli, I. Smarter, Better, Faster, Longer: A Modern Bidirectional Encoder for Fast, Memory Efficient, and Long Context Finetuning and Inference. \textit{arXiv preprint} arXiv:2412.13663, 2024.

\bibitem[B. Li et~al.(2024)]{inferenceserving}
Li, B., Jiang, Y., Gadepally, V., and Tiwari, D. LLM Inference Serving: Survey of Recent Advances and Opportunities. \textit{arXiv preprint} arXiv:2407.12391, 2024.

\bibitem[Trivedi et~al.(2022)]{trivedi2022musique}
Trivedi, H., Balasubramanian, N., Khot, T., and Sabharwal, A. MuSiQue: Multihop Questions via Single-hop Question Composition. \textit{Transactions of the ACL}, 10:539--554, 2022.

\bibitem[Gao et~al.(2025)]{kvcompresssurvey}
Gao, W., Zhou, X., Sun, P., Zhang, T., and Wen, Y. Rethinking Key-Value Cache Compression Techniques for Large Language Model Serving. \textit{arXiv preprint} arXiv:2503.24000, 2025.

\bibitem[X. Li et~al.(2025)]{kvtuner}
Li, X., Xing, Z., Li, Y., Qu, L., Zhen, H.-L., Liu, W., Yao, Y., Pan, S. J., and Yuan, M. KVTuner: Sensitivity-Aware Layer-Wise Mixed-Precision KV Cache Quantization for Efficient and Nearly Lossless LLM Inference. \textit{arXiv preprint} arXiv:2502.04420, 2025.

\bibitem[Cheng et~al.(2025)]{lmcache}
Cheng, Y., Liu, Y., Yao, J., An, Y., Chen, X., Feng, S., Huang, Y., Shen, S., Du, K., and Jiang, J. LMCache: An Efficient KV Cache Layer for Enterprise-Scale LLM Inference. \textit{arXiv preprint} arXiv:2510.09665, 2025.

\bibitem[Zhong et~al.(2024)]{zhong2024distserve}
Zhong, Y. et al. DistServe: Disaggregating Prefill and Decoding for Goodput-optimized Large Language Model Serving. In \textit{18th USENIX OSDI}, 2024.

\end{thebibliography}

\end{document}
